\chapter*{Preface}
\addcontentsline{toc}{chapter}{Preface}

Real model-based AI work begins with a problem. What decision are we supporting? What evidence do we have? What claim are we making? And what would justify trusting the result?

This book is organized around those questions. It is problem-driven rather than algorithm-driven. Models, losses, architectures, and systems appear because they resolve modeling dilemmas, and each method enters through the decision, data, representation, evidence, and system constraints that make it useful.

The book is written for undergraduate students, teachers, and self-learners who want a coherent first introduction to machine learning, evaluation, and reliable AI systems, and for IOAI students and early graduate readers who want the algorithmic and mathematical side of each method taken seriously. It assumes curiosity, some programming experience, and enough mathematical maturity to follow carefully explained equations. Two mathematical bridge chapters review linear algebra and probability, risk, and learning signals. Most first-course readers should begin with Chapter 1 and consult those bridges as just-in-time review; advanced readers can use them as quick reference.

A practical way to read the book. \textbf{For a first course}, take Chapters~\ref{ch:framing-learning-problems}--\ref{ch:dataset-design} in order, then Chapter~\ref{ch:optimization-neural-networks}, one modality chapter (Chapter~\ref{ch:vision}, \ref{ch:language-grounding}, or \ref{ch:tabular}), and a light pass through Chapter~\ref{ch:experiments-claims}. Use the two Mathematical Bridge chapters as just-in-time review when notation slows you down. \textbf{For a more advanced reading}---IOAI training, a graduate course, or a second pass---continue with Chapters~\ref{ch:probabilistic-models}, \ref{ch:representation-foundation-models}, \ref{ch:generative-representations}, \ref{ch:reinforcementlearning}, \ref{ch:audio-time}, \ref{ch:ranking-recommendation}, and~\ref{ch:reliability}, which carry the derivations, proofs, and from-scratch implementations. Chapters~\ref{ch:experiments-claims}--\ref{ch:models-systems} close both paths with experiments, reliability, and systems.

The center of gravity is modern machine learning: probabilistic models, generative models, representation learning, recommendation systems, reinforcement learning, and the system habits needed to evaluate, stress-test, and operationalize model-based systems.

Every serious AI project eventually has to answer seven questions:
\begin{enumerate}[leftmargin=*,label=\arabic*.]
\item What decision are we actually supporting?
\item What learning task is the closest defensible proxy for that decision?
\item What data stands in for the world, and what could bias or leak it?
\item What representation makes the relevant structure visible?
\item What is the simplest serious baseline?
\item What evidence would justify the claim?
\item How will the model output become an action inside a real system?
\end{enumerate}

Worked implementations for every chapter live in the companion repository.

\section*{About This Project}

This book began as lecture notes for IOAI students, undergraduates, and graduate students, and was later expanded into a textbook. It is published as an open-access educational resource through NOKI, the AI Olympiad in Norway, which trains the Norwegian national team for the International Olympiad in Artificial Intelligence (IOAI). NOKI is directed by the author, Sam Urmian, a researcher at the University of Bergen.

The book text is released under \textbf{CC~BY-NC-SA~4.0}; the companion code under \textbf{MIT}. The companion repository carries chapter-aligned implementations in three layers---minimal, practical, and extended---alongside an instructor guide with course paths and sample solutions.

\begin{flushleft}
Repository: \url{https://github.com/mlgorithm/ml-by-design}\\
Zenodo (concept DOI): \url{https://doi.org/10.5281/zenodo.19341954}\\
Contact: \texttt{sam.urmian@gmail.com}
\end{flushleft}
